\documentclass[UTF8]{ctexart}

\usepackage[a4paper,margin=1in]{geometry}
\usepackage{amsmath,amssymb,mathtools,bm}
\usepackage{booktabs,longtable,array,tabularx,makecell,multirow}
\usepackage{hyperref}
\usepackage{xurl}
\usepackage{enumitem}
\usepackage{float}

\hypersetup{
    colorlinks=true,
    linkcolor=black,
    urlcolor=blue,
    citecolor=black
}

\setlist[itemize]{leftmargin=2em}
\setlist[enumerate]{leftmargin=2em}
\setlength{\LTleft}{0pt}
\setlength{\LTright}{0pt}
\setlength{\emergencystretch}{3em}
\renewcommand{\arraystretch}{1.15}
\allowdisplaybreaks
\sloppy

\newcolumntype{Y}{>{\raggedright\arraybackslash}X}
\newcommand{\codecell}[1]{\parbox[t]{\linewidth}{\raggedright\ttfamily\small #1}}
\newcommand{\wrapcell}[1]{\parbox[t]{\linewidth}{\raggedright #1}}

\title{mra.py 中图学习插补模块的科研式说明手册}
\author{}
\date{2026-04-13}

\begin{document}

\maketitle

\begin{abstract}
本文面向仓库当前版本的 \texttt{mra.py}，系统说明其中与\textbf{图学习插补}直接相关的代码、数学原理、训练流程与测试流程。文档重点分析四条主线。第一，图插补模块的代码入口、训练入口与测试入口分别位于何处，真实执行调用链怎样串联。第二，图学习功能如何从训练窗口构造节点动态上下文、距离先验与自适应邻接矩阵，以及这张图如何驱动 GCN 完成窗口内变量关系建模。第三，训练阶段哪些参数会被更新、由哪些损失更新、这些训练得到的参数和统计量在测试集里怎样继续发挥作用。第四，代码中的张量在各子模块之间如何流动，输入输出形状如何变化。文中所有关键结论均以源代码与一次真实运行探针为依据，并额外标注当前实现中的重要细节，例如训练期额外遮挡、图先验的训练集来源、\texttt{detector\_loss} 对图分支的反向作用，以及 \texttt{node\_coords} 在当前主程序路径下实际上不参与前向计算这一事实。频域专家与检测器只在解释图模块耦合关系时作必要说明，不展开全部细节。
\end{abstract}

\noindent\textbf{关键词：} 图学习插补，GCN，自适应邻接矩阵，多采样率，门控融合，异常检测

\tableofcontents
\newpage

\section{文档范围与核心结论}

\subsection{分析范围}

本文只分析与图学习插补主线直接相关的内容：
\begin{itemize}
    \item \texttt{mra.py} 中 \texttt{PreImpute}、\texttt{TwoExpertGatedImputer}、\texttt{FusionAnomalyModel}、\texttt{compute\_losses}、\texttt{train\_model}、\texttt{reconstruct\_full\_sequence}、\texttt{collect\_window\_statistics}；
    \item \texttt{graph\_gcn\_reconstruction.py} 中 \texttt{ObservedStandardScaler}、训练集距离先验构造、\texttt{GraphLearner}、\texttt{GraphGCNReconstructor}、\texttt{adjacency\_balance\_loss}；
    \item \texttt{utils/methods/windowing.py} 与 \texttt{utils/methods/data\_loading.py} 中决定训练窗口、测试窗口和结构性缺失掩码含义的函数；
    \item 图专家与频域专家、检测器之间的耦合位置，但不会把频域专家和 Transformer 检测器展开成全文主体。
\end{itemize}

\subsection{先给结论}

基于代码阅读和真实运行探针，可以先给出九条最关键的结论：
\begin{itemize}
    \item 图插补模块的\textbf{代码入口}位于 \texttt{TwoExpertGatedImputer.forward} 中的
    \begin{quote}
        \ttfamily\small
        x\_gcn, adjacency = self.gcn(x\_seed, model\_missing\_mask)
    \end{quote}
    这里的 \texttt{self.gcn} 实际上是 \texttt{GraphGCNReconstructor}。
    \item 运行 \texttt{python mra.py} 时，真正发生参数更新的\textbf{训练入口}不是独立实验脚本，而是 \texttt{mra.py} 中的 \texttt{train\_model}；图专家、频域专家、门控网络和检测器共享一个 AdamW 优化器，属于\textbf{联合训练}。
    \item 图学习器不是在全局上学习一张固定邻接矩阵，而是\textbf{对每个窗口动态生成一张变量图}。这张图依赖于该窗口内的均值、标准差、最近一次观测值和缺失比例，因此是窗口级、自适应、按样本变化的。
    \item 在当前 \texttt{mra.py} 主程序路径中，图先验矩阵总会先由训练集构造出来，然后传给 \texttt{GraphLearner}。因此 \texttt{GraphLearner.node\_coords} 这组参数虽然被定义为可训练参数，但在当前主路径下并不参与前向，也不会获得梯度。
    \item 图学习分支的数学核心不是传统固定图卷积，而是“\textbf{训练集距离先验 + 窗口动态统计 + 节点相似度}”共同决定的邻接打分函数，再通过温度化 softmax、对称化和行归一化得到邻接矩阵。
    \item GCN 重构时使用的是\textbf{同一窗口共享、对时间维共享}的邻接矩阵：时间维没有在图卷积里被卷，图分支真正建模的是“同一时刻不同变量之间的关系”；时间信息主要在构图时通过窗口统计量进入。
    \item 五项损失对参数的作用范围不同：\texttt{gcn\_loss} 只更新图分支，\texttt{graph\_loss} 只更新图学习器，\texttt{fusion\_loss} 会同时更新图分支、频域分支、门控和采样率嵌入，\texttt{detector\_loss} 还会进一步反向更新图分支和门控。
    \item 训练集不仅更新神经网络权重，还生成一批\textbf{测试阶段必须继续使用}的非梯度量：标准化器均值方差、训练集驱动的距离先验矩阵、\texttt{rate\_id}、\texttt{stride}、训练分数均值与标准差、异常阈值。
    \item 当前默认数据下，训练集和测试集均为 \(4100 \times 18\)，插补窗口均为 \(4100 \times 50 \times 18\)，训练集中变量被自动识别为三类采样率：步长 \(1\)、\(3\)、\(5\)。
\end{itemize}

\section{代码入口、训练入口与总体调用链}

\subsection{入口总表}

{\small
\begin{longtable}{@{}>{\raggedright\arraybackslash}p{0.21\linewidth} >{\raggedright\arraybackslash}p{0.21\linewidth} >{\raggedright\arraybackslash}p{0.50\linewidth}@{}}
\toprule
位置 & 代码定位 & 作用 \\
\midrule
\endhead
主程序入口 & \codecell{mra.py\\:862--1122} & 读取训练/测试数据、标准化、构建窗口、推断采样率元数据、训练模型、补全序列、统计分数、阈值化并保存结果。 \\
采样率元数据入口 & \codecell{mra.py\\:150--161} & 根据训练集缺失率推断每个变量的 \texttt{stride} 和 \texttt{rate\_id}。 \\
训练期额外遮挡 & \codecell{mra.py\\:173--193} & 在可观测位置上随机采样 \texttt{holdout\_mask}，构造自监督插补目标。 \\
预插补入口 & \codecell{mra.py\\:196--259} & 通过双向填充生成 \(\bm{X}_{\text{seed}}\)，为图专家与频域专家提供平滑输入。 \\
图专家代码入口 & \codecell{mra.py\\:317--319} & 在 \texttt{TwoExpertGatedImputer.forward} 中调用 \texttt{GraphGCNReconstructor}。 \\
门控融合入口 & \codecell{mra.py\\:325--343} & 拼接图专家、频域专家与采样率嵌入，输出门控权重和最终插补结果。 \\
联合模型封装 & \codecell{mra.py\\:470--524} & 用 \texttt{FusionAnomalyModel} 串起插补器与检测器。 \\
损失定义入口 & \codecell{mra.py\\:527--560} & 定义 \texttt{fusion/gcn/freq/detector/graph} 五项损失和总损失。 \\
真实训练入口 & \codecell{mra.py\\:563--636} & 前向、损失、反向传播、梯度裁剪、优化器更新都在这里完成。 \\
测试补全入口 & \codecell{mra.py\\:639--677} & 用训练后的模型输出补全序列、门控权重和邻接矩阵样本。 \\
测试评分入口 & \codecell{mra.py\\:680--737} & 计算 \texttt{detector\_score}、\texttt{disagreement\_score} 和 \texttt{gate\_entropy}。 \\
训练集距离先验 & \codecell{graph\_gcn\_\\reconstruction.py\\:228--273} & 用训练数据轨迹和观测模式构造变量之间的欧氏距离矩阵。 \\
图学习定义 & \codecell{graph\_gcn\_\\reconstruction.py\\:350--458} & 定义 \texttt{GraphLearner}，根据窗口动态上下文生成邻接矩阵。 \\
GCN 重构定义 & \codecell{graph\_gcn\_\\reconstruction.py\\:463--506} & 定义 \texttt{GraphGCNReconstructor}，执行三层图消息传递重构。 \\
\bottomrule
\end{longtable}
}

\subsection{真实调用链}

当用户运行 \texttt{python mra.py} 时，与图插补最相关的真实执行链条为
\begin{equation}
    \begin{aligned}
        \texttt{main}
        &\rightarrow \texttt{train\_model}
        \rightarrow \texttt{FusionAnomalyModel.forward} \\
        &\rightarrow \texttt{TwoExpertGatedImputer.forward} \\
        &\rightarrow \texttt{PreImpute.fill} \\
        &\rightarrow \texttt{GraphGCNReconstructor.forward} \\
        &\rightarrow \texttt{GraphLearner.forward} \\
        &\rightarrow \texttt{compute\_losses}
        \rightarrow \texttt{loss.backward}
        \rightarrow \texttt{AdamW.step}.
    \end{aligned}
\end{equation}

这条链路有两个含义：
\begin{enumerate}
    \item 图专家不是独立于整套模型外部训练的，而是被包进联合模型内部；
    \item 图分支的训练结果不只影响插补误差，还会通过门控和检测器继续影响最终异常分数。
\end{enumerate}

\subsection{图专家与其它模块的边界}

虽然图学习是本文重点，但当前实现中的图专家并不单独闭环。它的输出 \(\bm{X}_{\text{gcn}}\) 至少会走向三条下游路径：
\begin{itemize}
    \item 与频域专家竞争门控权重，参与 \(\bm{X}_{\text{imputed}}\) 和 \(\bm{X}_{\text{complete}}\)；
    \item 与频域专家做差，构成 \texttt{disagreement\_score}；
    \item 通过 \(\bm{X}_{\text{complete}}\) 进入检测器，间接影响 \texttt{detector\_score} 和后续阈值化结果。
\end{itemize}

因此要真正理解图学习模块，必须把它放在“图专家 + 门控 + 检测器”的联合训练系统中分析，而不是只盯住 \texttt{GraphLearner} 自身。

\section{数据、窗口、掩码与训练测试切分}

\subsection{默认数据的真实尺寸}

对仓库当前默认数据，代码读取后得到：
\begin{align}
    \bm{X}^{\text{train}}_{\text{raw}} &\in \mathbb{R}^{4100 \times 18}, \\
    \bm{X}^{\text{test}}_{\text{raw}} &\in \mathbb{R}^{4100 \times 18}, \\
    \bm{M}^{\text{train}}_{s},\bm{M}^{\text{test}}_{s} &\in \{0,1\}^{4100 \times 18}.
\end{align}

其中 \(\bm{M}_{s}\) 是\textbf{结构性缺失掩码}，也就是 CSV 中原生的 \texttt{NaN} 位置。当前训练集和测试集的总体缺失率都约为 \(0.4889\)。

\subsection{只用观测值拟合的标准化器}

\texttt{ObservedStandardScaler} 只用训练集观测值拟合每个变量的均值 \(\mu_n\) 与标准差 \(\sigma_n\)，然后对训练集和测试集都做同样的变换：
\begin{equation}
    X_{t,n} =
    \begin{cases}
        \dfrac{X^{\text{raw}}_{t,n} - \mu_n}{\sigma_n}, & M_{s,t,n}=0, \\[6pt]
        0, & M_{s,t,n}=1.
    \end{cases}
\end{equation}

这意味着标准化后的数值零有两层可能语义：
\begin{itemize}
    \item 某个观测值在标准化后恰好等于零；
    \item 某个位置原本缺失，被训练集均值填充后标准化成零。
\end{itemize}

因此后续模型必须同时读取数值张量和掩码张量，不能只看数值本身。

\subsection{窗口构造方式}

当前代码中有三类窗口：
\begin{itemize}
    \item \texttt{build\_windows}：用于训练插补与完整序列补全，允许在序列前端重复首行进行补齐；
    \item \texttt{build\_standard\_windows}：用于训练阶段统计异常分数，不做前端补齐；
    \item \texttt{build\_prompt\_test\_windows}：用于测试阶段统计异常分数，同时生成固定标签。
\end{itemize}

默认 \texttt{seq\_len=50}、\texttt{stride=1} 时，真实形状为
\begin{align}
    \bm{W}^{\text{train}}_{\text{imp}} &\in \mathbb{R}^{4100 \times 50 \times 18}, \\
    \bm{W}^{\text{test}}_{\text{imp}} &\in \mathbb{R}^{4100 \times 50 \times 18}, \\
    \bm{W}^{\text{train}}_{\text{eval}} &\in \mathbb{R}^{4001 \times 50 \times 18}, \\
    \bm{W}^{\text{test}}_{\text{eval}} &\in \mathbb{R}^{4000 \times 50 \times 18}.
\end{align}

其中训练插补窗口与测试插补窗口都覆盖全序列，因此 \texttt{reconstruct\_full\_sequence} 只要取每个窗口的最后一个时间步，就能对齐还原出长度为 \(4100\) 的完整序列。

\subsection{训练集自动推断的多采样率元数据}

\texttt{infer\_rate\_metadata} 根据每列缺失率估计观测率
\begin{equation}
    r_n = 1 - \frac{1}{T_{\text{all}}}\sum_{t=1}^{T_{\text{all}}} M_{s,t,n},
\end{equation}
并据此得到采样步长
\begin{equation}
    s_n = \operatorname{round}\!\left(\frac{1}{\max(r_n,10^{-6})}\right).
\end{equation}

对当前默认数据，18 个变量被分为三组：

\begin{center}
\small
\begin{tabular}{cccccc}
\toprule
变量组 & 列索引 & 缺失率 & 观测率 & 推断步长 \(s_n\) & \texttt{rate\_id} \\
\midrule
高频组 & 1--6 & 0.0 & 1.0 & 1 & 0 \\
中频组 & 7--12 & 0.6667 & 0.3333 & 3 & 1 \\
低频组 & 13--18 & 0.8 & 0.2 & 5 & 2 \\
\bottomrule
\end{tabular}
\end{center}

因此主程序中会固定使用
\begin{equation}
    \bm{stride} = [1,1,1,1,1,1,3,3,3,3,3,3,5,5,5,5,5,5],
\end{equation}
\begin{equation}
    \bm{rate\_id} = [0,0,0,0,0,0,1,1,1,1,1,1,2,2,2,2,2,2].
\end{equation}

需要特别指出：图学习器本身\textbf{并不显式读取} \(\bm{rate\_id}\) 或 \(\bm{stride}\)。图分支只读取 \(\bm{X}_{\text{seed}}\) 与掩码，因此它对多采样率的感知是\textbf{隐式的}，主要来自缺失模式和窗口统计；显式采样率嵌入只出现在门控层与检测器中。

\section{图插补模块的总体结构与张量流}

\subsection{训练 batch 中的三种掩码}

对一个训练 batch，记真实标准化窗口为
\begin{equation}
    \bm{X}_{\text{true}} \in \mathbb{R}^{B \times T \times N},
\end{equation}
其中默认 \(T=50\)、\(N=18\)。代码中会同时出现三种掩码：
\begin{itemize}
    \item \(\bm{M}_{s}\)：结构性缺失掩码，来自原始数据；
    \item \(\bm{M}_{h}\)：训练期额外遮挡掩码，只从可观测位置里抽样；
    \item \(\bm{M}\)：模型输入掩码，满足 \(\bm{M}=\max(\bm{M}_{s},\bm{M}_{h})\)。
\end{itemize}

于是模型真正接收到的输入窗口为
\begin{equation}
    \bm{X}_{\text{input}} = \bm{X}_{\text{true}} \odot (1-\bm{M}).
\end{equation}

\texttt{sample\_holdout\_mask} 还包含一个非常重要的保底逻辑：如果某个样本窗口里本来有观测值，但按伯努利采样没有抽到任何 holdout 点，那么代码会\textbf{强行随机挑一个观测点做 holdout}。这保证了插补损失在该样本上不会完全失去监督。

\subsection{整体前向流程}

图插补与其上下游的关系可以写成
\begin{equation}
    \begin{aligned}
        \bm{X}_{\text{seed}} &= \text{PreImpute}(\bm{X}_{\text{input}},\bm{M}), \\
        \bm{X}_{\text{gcn}}, \bm{A} &= f_{\text{graph}}(\bm{X}_{\text{seed}},\bm{M}), \\
        \bm{X}_{\text{freq}} &= f_{\text{freq}}(\bm{X}_{\text{seed}}), \\
        \bm{\Pi} &= g(\bm{X}_{\text{seed}},\bm{M},\bm{X}_{\text{gcn}},\bm{X}_{\text{freq}},\bm{rate\_id}), \\
        \bm{X}_{\text{imputed}} &= \Pi_{:,:,:,0}\odot \bm{X}_{\text{gcn}} + \Pi_{:,:,:,1}\odot \bm{X}_{\text{freq}}, \\
        \bm{X}_{\text{complete}} &= \bm{M}\odot \bm{X}_{\text{imputed}} + (1-\bm{M})\odot \bm{X}_{\text{input}}.
    \end{aligned}
\end{equation}

从这个公式可以直接看出三点：
\begin{itemize}
    \item 图专家输出的是 \(\bm{X}_{\text{gcn}}\) 与邻接矩阵 \(\bm{A}\)；
    \item 图专家不是直接生成最终结果，而是与频域专家一起被门控层加权；
    \item 在非缺失位置，\(\bm{X}_{\text{complete}}\) 保持原始输入值；只有在结构性缺失或 holdout 位置，模型才会写入插补结果。
\end{itemize}

\subsection{为什么先做 PreImpute}

\texttt{PreImpute.fill} 使用双向时域填充生成 \(\bm{X}_{\text{seed}}\)。这一步的目的不是做最终插补，而是给图专家和频域专家提供一个更平滑、可处理的种子输入。它的规则是：
\begin{itemize}
    \item 若当前位置可观测，直接保留原值；
    \item 若缺失且前后两侧都有观测值，则取正向最近观测值与反向最近观测值的均值；
    \item 若仅一侧存在可用观测值，则取单侧填充值；
    \item 若整条可见历史都不存在，则保持零。
\end{itemize}

这一策略是\textbf{双向最近观测均值填充}，不是严格按时间位置做线性插值。举例而言，对单变量片段 \([1,\ ?,\ ?,\ 4]\)，该模块会得到 \([1,\ 2.5,\ 2.5,\ 4]\)，而不是 \([1,\ 2,\ 3,\ 4]\)。

\subsection{关键张量形状总表}

{\small
\begin{longtable}{@{}>{\raggedright\arraybackslash}p{0.23\linewidth} >{\raggedright\arraybackslash}p{0.22\linewidth} >{\raggedright\arraybackslash}p{0.47\linewidth}@{}}
\toprule
张量 & 默认形状 & 说明 \\
\midrule
\endhead
\(\bm{X}_{\text{true}}\) & \(\mathbb{R}^{B \times 50 \times 18}\) & 训练或测试窗口，值已按训练集统计量标准化。 \\
\(\bm{M}_{s},\bm{M}_{h},\bm{M}\) & \(\mathbb{R}^{B \times 50 \times 18}\) & 结构性缺失、额外遮挡、模型输入掩码。 \\
\(\bm{X}_{\text{input}}\) & \(\mathbb{R}^{B \times 50 \times 18}\) & 经过 \(\bm{M}\) 置零后的模型输入。 \\
\(\bm{X}_{\text{seed}}\) & \(\mathbb{R}^{B \times 50 \times 18}\) & 预插补结果，图专家与频域专家的共同输入。 \\
节点动态上下文 & \(\mathbb{R}^{B \times 18 \times 4}\) & 依次为均值、标准差、最近一次观测值、缺失比例。 \\
节点静态上下文 & \(\mathbb{R}^{B \times 18 \times 8}\) & 可训练的节点静态向量 \texttt{static\_context} 复制到 batch 维。 \\
节点表示 & \(\mathbb{R}^{B \times 18 \times 32}\) & \texttt{node\_encoder} 输出。 \\
边特征 & \(\mathbb{R}^{B \times 18 \times 18 \times 128}\) & 由 \([h_i,h_j,|h_i-h_j|,h_i \odot h_j]\) 拼接得到。 \\
邻接矩阵 \(\bm{A}\) & \(\mathbb{R}^{B \times 18 \times 18}\) & 每个窗口一张变量图。 \\
GCN 第一层输出 & \(\mathbb{R}^{B \times 50 \times 18 \times 32}\) & \texttt{gcn\_in} 后的隐藏表示。 \\
GCN 第二层输出 & \(\mathbb{R}^{B \times 50 \times 18 \times 32}\) & \texttt{gcn\_mid} 后的隐藏表示。 \\
\(\bm{X}_{\text{gcn}}\) & \(\mathbb{R}^{B \times 50 \times 18}\) & 图专家重构输出。 \\
\(\bm{X}_{\text{freq}}\) & \(\mathbb{R}^{B \times 50 \times 18}\) & 频域专家重构输出。 \\
门控输入 & \(\mathbb{R}^{B \times 50 \times 18 \times 14}\) & 6 个标量特征加 8 维采样率嵌入。 \\
门控权重 \(\bm{\Pi}\) & \(\mathbb{R}^{B \times 50 \times 18 \times 2}\) & 最后一维对应图专家和频域专家。 \\
\(\bm{X}_{\text{imputed}}\) & \(\mathbb{R}^{B \times 50 \times 18}\) & 两个专家的加权和。 \\
\(\bm{X}_{\text{complete}}\) & \(\mathbb{R}^{B \times 50 \times 18}\) & 用于下游检测器的完整窗口。 \\
\bottomrule
\end{longtable}
}

\section{图学习功能的数学原理}

\subsection{图先验首先来自训练集，而不是先验物理拓扑}

若用户没有提供外部物理坐标或距离矩阵，主程序会调用 \texttt{resolve\_distance\_prior\_matrix}，用训练集自动生成一个变量之间的距离矩阵。对第 \(n\) 个变量，代码构造描述向量
\begin{equation}
    \bm{u}_n =
    \left[
        \bm{x}^{\text{train}}_{:,n};
        \lambda \bm{o}^{\text{train}}_{:,n}
    \right]
    \in \mathbb{R}^{2T_{\text{all}}},
\end{equation}
其中 \(\bm{x}^{\text{train}}_{:,n}\) 是训练集标准化轨迹，\(\bm{o}^{\text{train}}_{:,n}=1-\bm{M}^{\text{train}}_{s,:,n}\) 是该变量的观测模式，默认 \(\lambda=0.5\)。

然后定义变量间距离
\begin{equation}
    D_{n,m} = \left\lVert \bm{u}_n - \bm{u}_m \right\rVert_2.
\end{equation}

这一步有两个非常重要的含义：
\begin{itemize}
    \item 默认图先验并不是管道布置、空间坐标或工艺流程图，而是\textbf{训练集驱动的变量相似性先验}；
    \item 该先验同时看轨迹相似度和缺失模式相似度，因此天然把多采样率信息的一部分折进了图先验。
\end{itemize}

\texttt{GraphLearner} 进一步把距离矩阵变成相似度先验
\begin{equation}
    P_{n,m} = \frac{1}{1 + D_{n,m}/\eta},
\end{equation}
其中 \(\eta\) 是距离矩阵正元素的中位数。由于对角线距离被强制置零，这个先验只作用于非对角项。

\subsection{窗口动态上下文的构造}

图学习器并不直接拿整个 \(50 \times 18\) 的窗口去生成边，而是先把每个变量在当前窗口中的信息压缩成 4 维动态统计量。对 batch 中第 \(b\) 个样本、变量 \(n\)，定义观测掩码
\begin{equation}
    O_{b,t,n} = 1 - M_{b,t,n}.
\end{equation}

再定义观测数
\begin{equation}
    c_{b,n} = \sum_{t=1}^{T} O_{b,t,n},
\end{equation}
并在代码中使用 \(\max(c_{b,n},1)\) 防止除零。随后可得：
\begin{align}
    \mu_{b,n} &= \frac{\sum_{t=1}^{T} X_{\text{seed},b,t,n} O_{b,t,n}}{\max(c_{b,n},1)}, \\
    \sigma_{b,n} &= \sqrt{
        \frac{\sum_{t=1}^{T}\left(X_{\text{seed},b,t,n}-\mu_{b,n}\right)^2 O_{b,t,n}}{\max(c_{b,n},1)}
        + 10^{-6}
    }, \\
    \ell_{b,n} &= \text{当前窗口中最近一次观测值}, \\
    \rho_{b,n} &= \frac{1}{T}\sum_{t=1}^{T} M_{b,t,n}.
\end{align}

于是节点动态上下文为
\begin{equation}
    \bm{d}_{b,n} = [\mu_{b,n},\sigma_{b,n},\ell_{b,n},\rho_{b,n}] \in \mathbb{R}^{4}.
\end{equation}

这个设计体现了当前图学习器的思想：\textbf{用窗口级统计量去概括变量在该窗口中的状态，再据此判断变量间连接强度}。也就是说，图不是由单个时间点的瞬时观测决定，而是由“最近一段时间内该变量的统计行为”决定。

\subsection{节点表示、边打分与自适应邻接矩阵}

除动态上下文外，\texttt{GraphLearner} 还为每个变量保存一个可训练静态向量 \(\bm{s}_n \in \mathbb{R}^{8}\)。于是节点输入为
\begin{equation}
    \bm{v}_{b,n} = [\bm{d}_{b,n};\bm{s}_n] \in \mathbb{R}^{12}.
\end{equation}

通过两层 MLP 得到节点表示
\begin{equation}
    \bm{h}_{b,n} = \phi_{\theta}(\bm{v}_{b,n}) \in \mathbb{R}^{32}.
\end{equation}

然后对任意节点对 \((n,m)\)，构造边特征
\begin{equation}
    \bm{p}_{b,n,m} =
    \left[
        \bm{h}_{b,n};
        \bm{h}_{b,m};
        |\bm{h}_{b,n}-\bm{h}_{b,m}|;
        \bm{h}_{b,n}\odot \bm{h}_{b,m}
    \right]
    \in \mathbb{R}^{128},
\end{equation}
并用 \texttt{edge\_mlp} 产生边打分
\begin{equation}
    E_{b,n,m} = \psi_{\theta}(\bm{p}_{b,n,m}).
\end{equation}

与此同时，代码还额外加入了节点表示的余弦相似度
\begin{equation}
    C_{b,n,m} =
    \frac{\bm{h}_{b,n}^{\top}\bm{h}_{b,m}}
    {\|\bm{h}_{b,n}\|_2 \|\bm{h}_{b,m}\|_2},
\end{equation}
并通过三个可训练标量控制先验、相似度和温度的强弱：
\begin{equation}
    \begin{aligned}
        \alpha &= \operatorname{softplus}(\texttt{prior\_strength}), \\
        \beta &= \operatorname{softplus}(\texttt{similarity\_strength}), \\
        \tau &= \operatorname{softplus}(\texttt{temperature}) + 10^{-4}.
    \end{aligned}
\end{equation}

令 \(\gamma = \operatorname{softplus}(\texttt{self\_loop\_logit})\)，则未归一化邻接 logit 可写成
\begin{equation}
    Z_{b,n,m} =
    \begin{cases}
        E_{b,n,m} + \alpha P_{n,m} + \beta C_{b,n,m}, & n \neq m, \\
        \gamma, & n = m.
    \end{cases}
\end{equation}

代码随后进行三步变换：
\begin{align}
    \tilde{A}_{b,n,m} &= \operatorname{softmax}_{m}\!\left(\frac{Z_{b,n,m}}{\tau}\right), \\
    \bar{A}_{b} &= \frac{1}{2}\left(\tilde{A}_{b} + \tilde{A}_{b}^{\top}\right), \\
    A_{b,n,m} &= \frac{\bar{A}_{b,n,m}}{\sum_{j}\bar{A}_{b,n,j}}.
\end{align}

\subsection{对上述构图公式的文字解释}

上面的公式虽然短，但设计含义很多，必须逐条解释：
\begin{enumerate}
    \item \textbf{边打分不是只靠一个来源。} 当前实现把边打分拆成三部分：数据驱动的边 MLP 输出 \(E_{b,n,m}\)、训练集距离先验 \(P_{n,m}\) 与窗口内节点表示余弦相似度 \(C_{b,n,m}\)。这让图既保留训练集结构偏好，又能根据当前窗口动态调整。
    \item \textbf{先验项是窗口不变的，动态项是窗口变化的。} \(P_{n,m}\) 对所有窗口共享，而 \(E_{b,n,m}\) 与 \(C_{b,n,m}\) 随当前窗口统计量变化。因此最终图是“固定先验 + 动态校正”的混合体。
    \item \textbf{softplus 使权重始终为正。} 这意味着代码不允许 \(\alpha\)、\(\beta\)、\(\tau\) 和 \(\gamma\) 取负数。直观地说，先验和相似度只能增强贡献强度，温度始终为正，自环 logit 始终为正。
    \item \textbf{softmax 让每个节点把权重分配给所有节点。} 行 softmax 后，每一行都变成非负且和为 1 的分布，因此每个节点都像是在决定“我应该从哪些变量吸收信息、各吸多少”。
    \item \textbf{对称化不是最终精确约束，而是中间修正。} 代码先把矩阵与其转置取平均，再做逐行归一化，因此最终输出矩阵是\textbf{近似对称且严格行归一化}的；最后一步后，严格意义上的完全对称一般不再精确成立。
    \item \textbf{图在时间维上是共享的。} 当前窗口里所有 50 个时间步共用同一张邻接矩阵。换句话说，图学习器认为“变量关系在这个窗口内近似稳定”，而不是每个时间步重新学一张图。
 \end{enumerate}

\subsection{图正则项的作用}

\texttt{adjacency\_balance\_loss} 只约束邻接矩阵对角元素接近目标值 \(\delta\)：
\begin{equation}
    \mathcal{L}_{\text{graph}} =
    \frac{1}{BN}\sum_{b=1}^{B}\sum_{n=1}^{N}\left(A_{b,n,n} - \delta\right)^2,
\end{equation}
默认 \(\delta=0.25\)。

由于每一行都归一化为 1，这个目标值的直观意义是：平均来说，希望每个节点保留约 \(25\%\) 的自环权重，把剩余 \(75\%\) 分配给其他变量。该正则并不直接鼓励稀疏性，也不直接约束拉普拉斯平滑性；它的主要作用是防止图退化成过强自环或过弱自环。

\section{GCN 重构与双专家门控的代码实现}

\subsection{GCN 重构的数学形式}

\texttt{GraphGCNReconstructor} 先用 \texttt{GraphLearner} 得到 \(\bm{A}\)，然后进行三层图消息传递。令
\begin{equation}
    \bm{H}^{(0)} = \bm{X}_{\text{seed}} \in \mathbb{R}^{B \times T \times N},
\end{equation}
并把其最后一维扩展成长度为 1 的特征维：
\begin{equation}
    \bm{H}^{(0)\prime} \in \mathbb{R}^{B \times T \times N \times 1}.
\end{equation}

\texttt{GCNLayer} 的计算方式是“先线性映射，再沿节点维用邻接矩阵聚合”。对任一层 \(l\)，其核心形式为
\begin{equation}
    \tilde{\bm{H}}^{(l)}_{b,t,n,:}
    =
    \sum_{m=1}^{N}
    A_{b,n,m}
    \left(
        \bm{H}^{(l-1)}_{b,t,m,:}\bm{W}^{(l)} + \bm{b}^{(l)}
    \right).
\end{equation}

第一层和第二层还会经过 \(\text{ReLU} + \text{LayerNorm} + \text{Dropout}\)：
\begin{align}
    \bm{H}^{(1)} &= \operatorname{Dropout}\!\left(
        \operatorname{LayerNorm}\!\left(
            \operatorname{ReLU}(\tilde{\bm{H}}^{(1)})
        \right)
    \right), \\
    \bm{H}^{(2)} &= \operatorname{Dropout}\!\left(
        \operatorname{LayerNorm}\!\left(
            \operatorname{ReLU}(\tilde{\bm{H}}^{(2)})
        \right)
    \right).
\end{align}

最后一层输出
\begin{equation}
    \bm{X}_{\text{gcn}} = \tilde{\bm{H}}^{(3)} \in \mathbb{R}^{B \times T \times N}.
\end{equation}

\subsection{对 GCN 设计的文字解释}

这里的 GCN 与很多图时序模型不同，重点有三条：
\begin{itemize}
    \item 图卷积只在\textbf{变量维}上传播，不在时间维上传播。因此它学习的是“同一时刻不同变量之间如何相互借力完成重构”。
    \item 时间信息没有完全消失，而是通过图学习器中的窗口统计量间接参与构图。也就是说，时间先用于\textbf{决定图}，图再用于\textbf{传播变量信息}。
    \item 由于当前 \(N=18\) 很小，显式构造 \(N \times N\) 的边特征和邻接矩阵计算开销是可接受的；如果变量数显著增加，该设计的二次复杂度会迅速抬升。
\end{itemize}

\subsection{门控层如何把图专家接入最终插补}

图专家并不是直接输出最终补全结果，而是与频域专家一起交给门控层决定权重。门控输入特征为
\begin{equation}
    \bm{g}_{b,t,n} =
    \left[
        X_{\text{seed},b,t,n},
        M_{b,t,n},
        X_{\text{gcn},b,t,n},
        X_{\text{freq},b,t,n},
        X_{\text{gcn},b,t,n} - X_{\text{seed},b,t,n},
        X_{\text{freq},b,t,n} - X_{\text{seed},b,t,n},
        \bm{e}^{\text{rate}}_{n}
    \right],
\end{equation}
其中 \(\bm{e}^{\text{rate}}_{n}\in\mathbb{R}^{8}\) 为采样率嵌入，因此默认总维度为 \(6+8=14\)。

经过两层 GELU MLP 后得到两位专家的 logit，再用 softmax 生成门控权重
\begin{equation}
    \bm{\pi}_{b,t,n} =
    \operatorname{softmax}\!\left(
        \text{MLP}_{\text{gate}}(\bm{g}_{b,t,n})
    \right)
    \in \mathbb{R}^{2}.
\end{equation}

于是
\begin{equation}
    X_{\text{imputed},b,t,n}
    =
    \pi_{b,t,n,0} X_{\text{gcn},b,t,n}
    +
    \pi_{b,t,n,1} X_{\text{freq},b,t,n}.
\end{equation}

最终完整窗口满足
\begin{equation}
    X_{\text{complete},b,t,n}
    =
    \begin{cases}
        X_{\text{input},b,t,n}, & M_{b,t,n}=0, \\
        X_{\text{imputed},b,t,n}, & M_{b,t,n}=1.
    \end{cases}
\end{equation}

这说明图专家对最终输出的影响并非“非黑即白”，而是由门控层按位置逐点决定的。

\section{训练流程、损失函数与参数更新}

\subsection{训练主循环}

\texttt{train\_model} 的每个 batch 都执行如下步骤：
\begin{enumerate}
    \item 从 \(\bm{W}^{\text{train}}_{\text{imp}}\) 取出 \(\bm{X}_{\text{true}}\) 和结构性缺失掩码 \(\bm{M}_{s}\)；
    \item 在观测位置上抽样 \(\bm{M}_{h}\)，形成 \(\bm{M}=\max(\bm{M}_{s},\bm{M}_{h})\)；
    \item 构造 \(\bm{X}_{\text{input}}=\bm{X}_{\text{true}}\odot(1-\bm{M})\)；
    \item 前向得到 \(\bm{X}_{\text{gcn}}\)、\(\bm{A}\)、\(\bm{X}_{\text{freq}}\)、\(\bm{X}_{\text{complete}}\) 和检测器重构结果；
    \item 计算五项损失并求和；
    \item \texttt{zero\_grad}，然后对总损失做 \texttt{backward}；
    \item 对全模型做 \(\ell_2\) 范数为 1 的梯度裁剪；
    \item 用 AdamW 更新参数。
\end{enumerate}

默认优化器为
\begin{equation}
    \text{AdamW}(\texttt{lr}=10^{-3},\ \texttt{weight\_decay}=10^{-4}).
\end{equation}

\subsection{五项损失的精确定义}

设 \(\bm{H}=\bm{M}_{h}\) 为 holdout 掩码，则三项插补监督损失均只在 holdout 位置上计算：
\begin{align}
    \mathcal{L}_{\text{fusion}}
    &= \operatorname{MSE}(\bm{X}_{\text{imputed}},\bm{X}_{\text{true}};\bm{H}), \\
    \mathcal{L}_{\text{gcn}}
    &= \operatorname{MSE}(\bm{X}_{\text{gcn}},\bm{X}_{\text{true}};\bm{H}), \\
    \mathcal{L}_{\text{freq}}
    &= \operatorname{MSE}(\bm{X}_{\text{freq}},\bm{X}_{\text{true}};\bm{H}).
\end{align}

图正则项为前节给出的
\begin{equation}
    \mathcal{L}_{\text{graph}} =
    \frac{1}{BN}\sum_{b,n}\left(A_{b,n,n}-\delta\right)^2.
\end{equation}

检测器损失定义为
\begin{equation}
    \mathcal{L}_{\text{det}}
    =
    \operatorname{MSE}
    \left(
        \bm{X}_{\text{det\_rec}},
        \operatorname{stopgrad}(\bm{X}_{\text{complete}});
        1-\bm{M}_{s}
    \right).
\end{equation}

这里有两个实现细节必须明确写出来：
\begin{itemize}
    \item 掩码用的是 \(1-\bm{M}_{s}\)，不是 \(1-\bm{M}\)。因此\textbf{训练期 holdout 位置仍然会参与检测器损失}；
    \item 目标张量 \(\bm{X}_{\text{complete}}\) 做了 \texttt{detach()}，但输入分支没有被截断，因此 \(\mathcal{L}_{\text{det}}\) 仍会反向更新插补器。
\end{itemize}

总代价函数为
\begin{equation}
    \mathcal{L}
    =
    1.0\,\mathcal{L}_{\text{fusion}}
    + 0.4\,\mathcal{L}_{\text{gcn}}
    + 0.4\,\mathcal{L}_{\text{freq}}
    + 0.5\,\mathcal{L}_{\text{det}}
    + 0.1\,\mathcal{L}_{\text{graph}}.
\end{equation}

\subsection{哪些损失更新哪些参数}

对当前代码做单 batch 梯度探针后，可以得到清晰的更新范围：

\begin{center}
\small
\begin{tabular}{lcccccc}
\toprule
损失项 & 图学习器 & GCN 层 & 频域专家 & 门控 MLP & 采样率嵌入 & 检测器 \\
\midrule
\(\mathcal{L}_{\text{fusion}}\) & 是 & 是 & 是 & 是 & 是 & 否 \\
\(\mathcal{L}_{\text{gcn}}\) & 是 & 是 & 否 & 否 & 否 & 否 \\
\(\mathcal{L}_{\text{freq}}\) & 否 & 否 & 是 & 否 & 否 & 否 \\
\(\mathcal{L}_{\text{det}}\) & 是 & 是 & 是 & 是 & 是 & 是 \\
\(\mathcal{L}_{\text{graph}}\) & 是 & 否 & 否 & 否 & 否 & 否 \\
\bottomrule
\end{tabular}
\end{center}

这张表意味着：
\begin{itemize}
    \item 图专家至少受到三种力共同作用：\(\mathcal{L}_{\text{gcn}}\) 的直接重构监督、\(\mathcal{L}_{\text{fusion}}\) 的门控融合监督，以及 \(\mathcal{L}_{\text{det}}\) 通过完整窗口传来的下游可重构性压力；
    \item \(\mathcal{L}_{\text{graph}}\) 不会直接训练 GCN 权重，它只训练 \texttt{GraphLearner} 自身输出更符合对角占比目标的图；
    \item 门控层不会被 \(\mathcal{L}_{\text{gcn}}\) 或 \(\mathcal{L}_{\text{freq}}\) 单独训练，它主要由 \(\mathcal{L}_{\text{fusion}}\) 和 \(\mathcal{L}_{\text{det}}\) 决定。
\end{itemize}

\subsection{图模块相关参数清单}

当前默认配置下，全模型共有 \(693{,}260\) 个参数，其中与图插补直接相关的主要参数如下：

\begin{center}
\small
\begin{tabular}{lccp{0.38\linewidth}}
\toprule
参数组 & 参数量 & 当前主路径是否实际更新 & 作用 \\
\midrule
\texttt{GraphLearner} & 5925 & 是，但 \texttt{node\_coords} 例外 & 生成窗口级邻接矩阵。 \\
\texttt{GCN} 三层与两层归一化 & 1281 & 是 & 按邻接矩阵在变量维上传播信息并输出 \(\bm{X}_{\text{gcn}}\)。 \\
门控 MLP & 5250 & 是 & 决定图专家与频域专家在每个位置的权重。 \\
插补器采样率嵌入 & 24 & 是 & 给门控层提供显式采样率信息。 \\
频域专家 & 24314 & 是 & 与图专家共同完成插补。 \\
检测器 & 656466 & 是 & 对 \(\bm{X}_{\text{complete}}\) 做重构并产生异常分数。 \\
\bottomrule
\end{tabular}
\end{center}

\texttt{GraphLearner} 内部又可以细分为：
\begin{itemize}
    \item 标量门控参数：\texttt{self\_loop\_logit}、\texttt{prior\_strength}、\texttt{similarity\_strength}、\texttt{temperature}；
    \item 节点静态上下文：\texttt{static\_context}；
    \item 节点表示网络：\texttt{node\_encoder}；
    \item 边打分网络：\texttt{edge\_mlp}；
    \item \texttt{node\_coords}。
\end{itemize}

其中 \texttt{node\_coords} 有 144 个参数，但在当前主程序路径中不会更新。原因很直接：\texttt{main} 总会先调用 \texttt{resolve\_distance\_prior\_matrix} 得到一个非空距离矩阵，并把它传给 \texttt{GraphLearner}；此时 \texttt{\_build\_prior\_adjacency} 只使用该固定矩阵，不再读取 \texttt{node\_coords}。

\section{训练集得到的参数和统计量在测试集中的作用}

\subsection{需要区分“梯度学习到的参数”和“训练集估计出的统计量”}

测试阶段真正起作用的量可以分成两类：
\begin{enumerate}
    \item \textbf{通过反向传播学习到的参数}：图学习器、GCN、门控层、频域专家、检测器的权重；
    \item \textbf{仅由训练集估计得到的固定量}：标准化均值方差、距离先验矩阵、\texttt{rate\_id}、\texttt{stride}、训练分数均值方差、异常阈值。
\end{enumerate}

把这两类混为一谈会直接导致对测试流程的理解错误。

\subsection{训练集到测试集的作用映射}

{\small
\begin{longtable}{@{}>{\raggedright\arraybackslash}p{0.24\linewidth} >{\raggedright\arraybackslash}p{0.24\linewidth} >{\raggedright\arraybackslash}p{0.44\linewidth}@{}}
\toprule
训练期得到的量 & 来源 & 在测试集中的作用 \\
\midrule
\endhead
\(\mu_n,\sigma_n\) & \codecell{ObservedStandardScaler.\\fit(train\_raw, train\_mask)} & 用同一套训练集统计量标准化测试集，并在输出 \texttt{test\_completed.csv} 前做逆变换。 \\
训练集驱动距离矩阵 \(\bm{D}\) & \codecell{resolve\_distance\_\\prior\_matrix(\\train\_scaled,\\train\_mask)} & 作为图先验输入 \texttt{GraphLearner}，影响测试窗口中每一张邻接矩阵的生成。 \\
\(\bm{rate\_id},\bm{stride}\) & \codecell{infer\_rate\_metadata(\\train\_mask)} & 在测试阶段继续作为门控层采样率嵌入和检测器相位编码的依据。 \\
图学习器、GCN、门控、频域专家、检测器参数 & \texttt{train\_model} 中 AdamW 训练得到 & 决定测试阶段的 \(\bm{X}_{\text{gcn}}\)、门控权重、\(\bm{X}_{\text{complete}}\)、检测器重构和异常分数。 \\
\texttt{train\_completed\_scaled} 及其窗口特征统计 & 训练集补全结果再窗口化后计算 & 用于构造 \texttt{shift\_score} 的训练参考分布，并把同一套统计量应用到测试补全窗口。 \\
训练分数均值方差 & \codecell{normalize\_scores(\\train, test)} & 把测试分数变成以训练集为基准的标准化分数。 \\
异常阈值 & \codecell{choose\_threshold(\\train\_scores, ...)} & 对测试最终分数做二值化，得到异常判定。 \\
\bottomrule
\end{longtable}
}

\subsection{测试阶段图模块怎样影响最终结果}

测试阶段虽然不再更新参数，但图模块仍通过三条路径持续发挥作用：
\begin{itemize}
    \item \textbf{路径 1：} 通过 \(\bm{X}_{\text{gcn}}\) 参与门控融合，影响 \(\bm{X}_{\text{complete}}\)，进而影响检测器重构误差和 \texttt{shift\_score}；
    \item \textbf{路径 2：} 直接与频域专家做差，形成 \texttt{disagreement\_score}；
    \item \textbf{路径 3：} 通过测试窗口动态生成邻接矩阵，改变图消息传递方式，从而影响补全结果和门控权重分布。
\end{itemize}

从最终打分公式看，
\begin{equation}
    s_{\text{final}}
    =
    \left|z_{\text{det}}\right|
    + \lambda_{1} z_{\text{gap}}
    + \lambda_{2} z_{\text{shift}},
\end{equation}
其中 \(z_{\text{det}}\)、\(z_{\text{gap}}\) 和 \(z_{\text{shift}}\) 分别表示以训练集均值和标准差归一化后的
\texttt{detector\_score}、\texttt{disagreement\_score} 和 \texttt{shift\_score}。
图模块至少影响前两项，并通过补全结果进一步影响第三项。

\section{测试流程、分数构造与图学习结果的输出}

\subsection{测试补全流程}

\texttt{reconstruct\_full\_sequence} 在测试阶段做的事情是：
\begin{enumerate}
    \item 不再使用随机 holdout，只把结构性缺失位置置零；
    \item 前向运行整套模型；
    \item 取每个窗口最后一个时间步的 \(\bm{X}_{\text{complete}}\)，拼成完整序列；
    \item 同时保存每个窗口最后一个时间步的门控权重，以及整张邻接矩阵样本。
\end{enumerate}

因此输出文件中：
\begin{itemize}
    \item \texttt{test\_completed.csv} 反映的是图专家与频域专家在真实测试缺失模式下的联合补全结果；
    \item \texttt{train\_gate\_weights.csv} 与 \texttt{test\_gate\_weights.csv} 反映图专家在不同变量上的平均权重；
    \item \texttt{train\_adjacency\_heatmap.png} 反映训练窗口邻接矩阵的平均图结构。
\end{itemize}

\subsection{测试评分流程}

\texttt{collect\_window\_statistics} 会为训练评分窗口和测试评分窗口都计算三类统计量：
\begin{align}
    \texttt{detector\_score}
    &= \frac{\sum (\bm{X}_{\text{det\_rec}}-\bm{X}_{\text{complete}})^2 \odot (1-\bm{M}_{s})}
    {\sum (1-\bm{M}_{s})}, \\
    \texttt{disagreement\_score}
    &= \operatorname{mean}\left|\bm{X}_{\text{gcn}}-\bm{X}_{\text{freq}}\right|, \\
    \texttt{gate\_entropy}
    &= -\operatorname{mean}\sum_{k=1}^{2}\Pi_{k}\log \Pi_{k}.
\end{align}

其中 \texttt{disagreement\_score} 是最直接体现图专家作用的测试统计量之一：如果图专家与频域专家在某些窗口上出现持续高分歧，那么这本身会被纳入最终异常分数。

\subsection{为什么训练集补全结果还会影响测试打分}

主程序在模型训练完毕后，还会对训练集和测试集分别做一次完整补全，再将这些补全后的窗口送入 \texttt{compute\_distribution\_shift\_scores}。也就是说，\texttt{shift\_score} 不是从原始带缺失窗口上计算出来的，而是从\textbf{补全后的窗口分布}上计算出来的。

于是图模块对测试 \texttt{shift\_score} 的影响来自两个方向：
\begin{itemize}
    \item 它决定训练补全窗口的参考分布是什么；
    \item 它决定测试补全窗口相对该参考分布偏移多大。
\end{itemize}

这也解释了为什么图插补模块虽然名义上是“插补前端”，却会显著影响最终异常检测性能。

\section{实现层面的设计思路、优点与局限}

\subsection{设计思路总结}

从代码实现看，图学习部分的设计思路可以概括成一句话：
\begin{quote}
    用训练集驱动的变量关系先验稳定图结构，用当前窗口的统计行为动态修正图结构，再用该图在变量维上传播信息完成插补。
\end{quote}

这套设计相比“固定邻接 + 固定 GCN”的优势主要有三点：
\begin{itemize}
    \item 图能够随窗口动态变化，更适合异常工况下关系漂移的场景；
    \item 先验矩阵提供了稳定的结构偏置，降低纯数据驱动构图的随机性；
    \item 图专家与频域专家通过门控互补，使图学习不必独自承担所有插补难点。
\end{itemize}

\subsection{当前实现的四个重要细节}

结合代码，可以额外指出四个容易被忽略但很重要的实现事实：
\begin{enumerate}
    \item \textbf{图先验默认来自训练数据而非物理拓扑。} 只有用户显式提供物理坐标或距离矩阵时，先验才会切换成外部物理先验。
    \item \textbf{图是窗口级的，不是时间步级的。} 这会降低复杂度，但也意味着窗口内部若变量关系快速变化，单张图可能不够细。
    \item \textbf{图正则只约束对角占比。} 当前没有显式稀疏约束，也没有图拉普拉斯正则，因此图结构的其它性质主要由数据误差和 softmax 温度间接塑造。
    \item \textbf{\texttt{node\_coords} 在当前主路径下是失活参数。} 这不是概念上的设计缺陷，但在实现上确实意味着有一组定义出的参数实际上没有被训练。
\end{enumerate}

\subsection{从研究复现角度给出的建议}

如果后续需要把该图模块发展为更完整的科研实现，优先可考虑以下方向：
\begin{itemize}
    \item 给 \texttt{GraphLearner} 增加显式稀疏约束或 top-\(k\) 边保留机制，提升图的可解释性；
    \item 若希望 \texttt{node\_coords} 发挥作用，可在训练集自动先验存在时把它作为可学习残差项加入先验，而不是完全绕开；
    \item 若窗口内关系变化很快，可把“按窗口一张图”改成“按时间步一张图”或“按若干子片段一张图”；
    \item 若想让图分支显式利用采样率，可把 \(\bm{rate\_id}\) 或相位编码并入节点动态上下文，而不仅仅交给门控层和检测器。
\end{itemize}

\section{结论}

回到用户最关心的问题，当前 \texttt{mra.py} 里的图插补模块可以被准确概括为：它先在训练集上拟合标准化器、采样率元数据和训练数据驱动的变量距离先验，再在每个窗口内根据均值、标准差、最近观测值和缺失比例生成一张自适应变量图，随后用这张图驱动三层 GCN 重构，并与频域专家通过门控融合，最终把补全结果送入检测器形成异常分数。

训练阶段更新的不只是图学习器和 GCN 参数，还包括门控层、频域专家和检测器；反过来，检测器损失也会回传影响图专家。测试阶段不再更新参数，但训练阶段得到的标准化统计量、图先验、采样率元数据、训练分数归一化参数和阈值都会继续生效。换言之，图学习模块既是前端插补器，也是后续异常检测分数形成过程中的结构性参与者，而不是一个与检测阶段割裂的独立预处理器。

\end{document}
